\documentclass[12pt]{article}

\input{~/School/Notes/header.tex}

\usepackage{booktabs}

\setlength{\parindent}{1.25cm}
\setlength{\parskip}{0.25em}
\renewcommand{\arraystretch}{1.15}

\title{Лабораторная работа по физике}
\author{Денисов Илья}
\date{}

\begin{document}
    
    \maketitle
    
    \vspace{5cm}
    {\centering \LARGE Моделирование случайной величины \\ и исследование её распределения \par}

    \newpage
    
    \tableofcontents
	\setcounter{tocdepth}{3}
	
	\newpage
    
    \section{Цель работы}
    
    Провести многократное измерение заданного промежутка времени, получить случайную выборку, найти её основные статистические характеристики, построить гистограмму плотности относительной частоты и сравнить полученное распределение с нормальным распределением.
    
    \section{Задачи}
    
    \begin{enumerate}
        \item Выполнить 100 измерений заданного промежутка времени.
        \item Найти среднее значение $\bar t$, среднеквадратичное отклонение $\sigma$ и среднеквадратичную погрешность среднего $\sigma_{\bar t}$.
        \item Построить гистограмму плотности относительной частоты результатов измерений.
        \item На тот же график нанести нормальное распределение с теми же $\bar t$ и $\sigma$.
        \item Проверить попадание результатов в интервалы $\bar t \pm \sigma$, $\bar t \pm 2 \sigma$ и $\bar t \pm 3 \sigma$.
        \item Найти случайную и полную погрешности измерения для нескольких доверительных вероятностей.
    \end{enumerate}
    
    \section{Метод измерения и оборудование}
    
    В работе использовались два таймера с погрешностью $0.01$ с.
    
    Первый таймер оставлялся запущенным. Затем по команде <<старт>> запускался второй таймер. Человек, контролировавший первый таймер, через заданные $5$ с подавал команду <<стоп>>, после чего второй таймер останавливался. Показание второго таймера записывалось в таблицу. Перед каждым следующим измерением обнулялся только второй таймер.
    
    Всего выполнено $N = 100$ измерений.
    
    \section{Рабочие формулы}
    
    Среднее значение измеренной величины:
    \[
    	\bar t = \frac{1}{N} \sum_{i = 1}^{N} t_i.
    \]
    
    Среднеквадратичное отклонение выборки:
    \[
    	\sigma = \sqrt{\frac{\displaystyle\sum_{i = 1}^{N}(t_i - \bar t)^2}{N - 1}}.
    \]
    
    Среднеквадратичная погрешность среднего значения:
    \[
    	\sigma_{\bar t} = \frac{\sigma}{\sqrt N}.
    \]
    
    Для нормального распределения используется функция
    \[
    	\rho(t) = \frac{1}{\sqrt{2\pi}\sigma} e^{-\dfrac{(t - \bar t)^2}{2 \sigma^2}}.
    \]
    
    Максимальное значение нормальной плотности достигается при $t = \bar t$:
    \[
    	\rho_{\max} = \frac{1}{\sqrt{2 \pi} \sigma}.
    \]
    
    Для гистограммы плотность относительной частоты в $i$-м интервале вычисляется по формуле
    \[
    	\rho_i^{\mathrm{hist}} = \frac{N_i}{N \, \Delta t},
    \]
    где $N_i$~--- число измерений, попавших в интервал ширины $\Delta t$.
    
    Случайная абсолютная погрешность при доверительной вероятности $\alpha$:
    \[
    	\Delta t_{\text{сл}} = t_{\alpha, N} \, \sigma_{\bar t},
    \]
    где $t_{\alpha, N}$~--- коэффициент Стьюдента.
    
    Относительные погрешности и полная погрешность:
    \[
    	\varepsilon_{\text{сл}} = \frac{\Delta t_{\text{сл}}}{\bar t}, \qquad
    	\varepsilon_{\text{пр}} = \frac{\Delta t_{\text{пр}}}{\bar t}, \qquad
    	\varepsilon_{\text{полн}} = \sqrt{\varepsilon_{\text{сл}}^2 + \varepsilon_{\text{пр}}^2}, \qquad
    	\Delta t_{\text{полн}} = \varepsilon_{\text{полн}} \, \bar t.
    \]
    где приборная погрешность $\Delta t_{\text{пр}} = 0.01$ с.
    
    \section{Результаты прямых измерений и их обработка}

	\begin{table}[H]
    	\centering
    	\begin{xltabular}{\textwidth}{@{}Xr@{}}
        	\toprule
        	Величина & Значение\\
        	\midrule
        	Число измерений $N$ & 100\\
        	Минимальное значение $t_{\min}$ & 4.49 с\\
        	Максимальное значение $t_{\max}$ & 5.40 с\\
        	Среднее $\bar t$ & 4.943 с\\
        	Дисперсия $D = \sigma^2$ & 0.026853 с$^2$\\
        	СКО $\sigma$ & 0.163869 с\\
        	СКО среднего $\sigma_{\bar t}$ & 0.016387 с\\
        	\bottomrule
    	\end{xltabular}
		\caption{Основные результаты обработки выборки}
    	\label{tab:main}
	\end{table}

	Среднее значение выборки равно
	\[
		\bar t = 4.9429 \ \text{с} \approx 4.943 \ \text{с}.
	\]

	Для таблицы отклонений используется округлённое среднее $4.943$ с. Тогда
	\[
		\sum_{i = 1}^{N} (t_i - \bar t) = -10 \ \text{мс}.
	\]
	По модулю это меньше допустимого значения $25$ мс, указанного в методических указаниях как контроль отклонения, связанного с округлением среднего.

	\subsection{Полная таблица 100 измерений}

	\small
	\begin{xltabular}{\textwidth}{@{}X X X X@{}}
		\toprule
		$N$ & $t_i$, с & $t_i - \bar t$, мс & $(t_i - \bar t)^2$, мс$^2$\\
		\midrule
		\endfirsthead
		\toprule
		$N$ & $t_i$, с & $t_i - \bar t$, мс & $(t_i - \bar t)^2$, мс$^2$\\
		\midrule
		\endhead
		\midrule
		\multicolumn{4}{r}{\textit{Продолжение на следующей странице}}\\
		\endfoot
		\midrule
		\multicolumn{3}{r}{\textbf{Контрольная сумма} $\sum (t_i - \bar t)$} & $-10$ мс \\
		\multicolumn{3}{r}{\textbf{Сумма квадратов}} & $2658460$ мс$^2$ \\
		\bottomrule
		\caption{Результаты измерений и отклонения от среднего}\label{tab:measurements}\\
		\endlastfoot
		1 & 5.34 & +397 & 157609 \\
		2 & 4.49 & -453 & 205209 \\
		3 & 4.97 & +27 & 729 \\
		4 & 5.19 & +247 & 61009 \\
		5 & 5.12 & +177 & 31329 \\
		6 & 4.95 & +7 & 49 \\
		7 & 5.40 & +457 & 208849 \\
		8 & 5.09 & +147 & 21609 \\
		9 & 5.02 & +77 & 5929 \\
		10 & 4.60 & -343 & 117649 \\
		11 & 5.10 & +157 & 24649 \\
		12 & 4.76 & -183 & 33489 \\
		13 & 4.82 & -123 & 15129 \\
		14 & 5.17 & +227 & 51529 \\
		15 & 4.91 & -33 & 1089 \\
		16 & 5.09 & +147 & 21609 \\
		17 & 4.75 & -193 & 37249 \\
		18 & 4.96 & +17 & 289 \\
		19 & 5.04 & +97 & 9409 \\
		20 & 5.12 & +177 & 31329 \\
		21 & 4.91 & -33 & 1089 \\
		22 & 5.22 & +277 & 76729 \\
		23 & 5.09 & +147 & 21609 \\
		24 & 4.99 & +47 & 2209 \\
		25 & 4.65 & -293 & 85849 \\
		26 & 5.23 & +287 & 82369 \\
		27 & 4.92 & -23 & 529 \\
		28 & 5.02 & +77 & 5929 \\
		29 & 5.17 & +227 & 51529 \\
		30 & 5.01 & +67 & 4489 \\
		31 & 4.85 & -93 & 8649 \\
		32 & 5.02 & +77 & 5929 \\
		33 & 5.04 & +97 & 9409 \\
		34 & 5.05 & +107 & 11449 \\
		35 & 4.80 & -143 & 20449 \\
		36 & 5.17 & +227 & 51529 \\
		37 & 4.82 & -123 & 15129 \\
		38 & 4.94 & -3 & 9 \\
		39 & 4.84 & -103 & 10609 \\
		40 & 4.87 & -73 & 5329 \\
		41 & 5.12 & +177 & 31329 \\
		42 & 4.69 & -253 & 64009 \\
		43 & 4.85 & -93 & 8649 \\
		44 & 4.84 & -103 & 10609 \\
		45 & 4.73 & -213 & 45369 \\
		46 & 4.72 & -223 & 49729 \\
		47 & 4.90 & -43 & 1849 \\
		48 & 4.82 & -123 & 15129 \\
		49 & 4.97 & +27 & 729 \\
		50 & 5.07 & +127 & 16129 \\
		51 & 4.77 & -173 & 29929 \\
		52 & 4.80 & -143 & 20449 \\
		53 & 4.97 & +27 & 729 \\
		54 & 4.74 & -203 & 41209 \\
		55 & 4.90 & -43 & 1849 \\
		56 & 4.84 & -103 & 10609 \\
		57 & 4.80 & -143 & 20449 \\
		58 & 5.17 & +227 & 51529 \\
		59 & 4.69 & -253 & 64009 \\
		60 & 4.90 & -43 & 1849 \\
		61 & 4.75 & -193 & 37249 \\
		62 & 4.92 & -23 & 529 \\
		63 & 5.04 & +97 & 9409 \\
		64 & 5.10 & +157 & 24649 \\
		65 & 5.10 & +157 & 24649 \\
		66 & 4.85 & -93 & 8649 \\
		67 & 4.94 & -3 & 9 \\
		68 & 4.75 & -193 & 37249 \\
		69 & 5.10 & +157 & 24649 \\
		70 & 4.92 & -23 & 529 \\
		71 & 4.94 & -3 & 9 \\
		72 & 5.02 & +77 & 5929 \\
		73 & 4.75 & -193 & 37249 \\
		74 & 5.07 & +127 & 16129 \\
		75 & 5.09 & +147 & 21609 \\
		76 & 4.70 & -243 & 59049 \\
		77 & 5.00 & +57 & 3249 \\
		78 & 4.72 & -223 & 49729 \\
		79 & 4.82 & -123 & 15129 \\
		80 & 4.87 & -73 & 5329 \\
		81 & 4.80 & -143 & 20449 \\
		82 & 4.85 & -93 & 8649 \\
		83 & 4.97 & +27 & 729 \\
		84 & 4.92 & -23 & 529 \\
		85 & 5.10 & +157 & 24649 \\
		86 & 4.99 & +47 & 2209 \\
		87 & 5.04 & +97 & 9409 \\
		88 & 4.90 & -43 & 1849 \\
		89 & 4.79 & -153 & 23409 \\
		90 & 4.97 & +27 & 729 \\
		91 & 4.74 & -203 & 41209 \\
		92 & 4.97 & +27 & 729 \\
		93 & 5.04 & +97 & 9409 \\
		94 & 5.20 & +257 & 66049 \\
		95 & 4.77 & -173 & 29929 \\
		96 & 5.05 & +107 & 11449 \\
		97 & 4.89 & -53 & 2809 \\
		98 & 5.10 & +157 & 24649 \\
		99 & 5.02 & +77 & 5929 \\
		100 & 4.94 & -3 & 9 \\
	\end{xltabular}
	\normalsize

	Отсюда
	\[
		\sigma = \sqrt{\frac{2658460 \ \text{мс}^2}{99}}
		\approx 0.163869 \ \text{с}.
	\]

	\section{Гистограмма и нормальное распределение}

	Ширина интервала гистограммы оценивается по формуле
	\[
		\Delta t = \frac{t_{\max} - t_{\min}}{\sqrt N} = \frac{5.40 - 4.49}{\sqrt{100}} = 0.091 \ \text{с}.
	\]
	Для удобства округлим её до
	\[
		\Delta t = 0.10 \ \text{с},
	\]
	а диапазон измерений $4.49 \dots 5.40$ с расширим до $4.45 \dots 5.45$ с, чтобы все результаты попали в $10$ интервалов вида $[t_k; \, t_k + \Delta t)$ с <<круглыми>> границами.

	\begin{table}[H]
		\centering
		\small
		\begin{xltabular}{\textwidth}{@{}X X X X X@{}}
			\toprule
			Интервал, с & $N_i$ & $\rho_i^{\mathrm{hist}}$, с$^{-1}$ &
			$\rho(a_i)$, с$^{-1}$ & $\rho(a_{i + 1})$, с$^{-1}$\\
			\midrule
			$[4.45; \, 4.55)$ & 1  & 0.10 & 0.026 & 0.137 \\
			$[4.55; \, 4.65)$ & 1  & 0.10 & 0.137 & 0.493 \\
			$[4.65; \, 4.75)$ & 9  & 0.90 & 0.493 & 1.218 \\
			$[4.75; \, 4.85)$ & 23 & 2.30 & 1.218 & 2.073 \\
			$[4.85; \, 4.95)$ & 17 & 1.70 & 2.073 & 2.432 \\
			$[4.95; \, 5.05)$ & 24 & 2.40 & 2.432 & 1.966 \\
			$[5.05; \, 5.15)$ & 15 & 1.50 & 1.966 & 1.095 \\
			$[5.15; \, 5.25)$ & 8  & 0.80 & 1.095 & 0.421 \\
			$[5.25; \, 5.35)$ & 1  & 0.10 & 0.421 & 0.111 \\
			$[5.35; \, 5.45]$ & 1  & 0.10 & 0.111 & 0.020 \\
			\bottomrule
		\end{xltabular}
		\caption{Данные для построения гистограммы и значения нормальной плотности на границах интервалов}
		\label{tab:hist}
	\end{table}

	\begin{figure}[H]
		\centering
		\begin{tikzpicture}
			\begin{axis}[
				width=0.9\textwidth,
				height=8cm,
				xlabel={$t$, с},
				ylabel={Число попаданий в интервал $N_i$},
				xmin=4.45, xmax=5.45,
				ymin=0, ymax=26.5,
				xtick={4.45,4.55,4.65,4.75,4.85,4.95,5.05,5.15,5.25,5.35,5.45},
				x tick label style={rotate=45, anchor=north east, font=\small},
				axis y line*=left,
				grid=major,
				]
				\addplot[
					ybar interval,
					draw=black,
					fill=gray!30,
				] coordinates {
					(4.45,1)
					(4.55,1)
					(4.65,9)
					(4.75,23)
					(4.85,17)
					(4.95,24)
					(5.05,15)
					(5.15,8)
					(5.25,1)
					(5.35,1)
					(5.45,1)
				};
				\label{plot:Ni}
			\end{axis}

			\begin{axis}[
				width=0.9\textwidth,
				height=8cm,
				xmin=4.45, xmax=5.45,
				ymin=0, ymax=2.65,
				axis y line*=right,
				axis x line=none,
				ylabel={Плотность относительной частоты $\rho_i$, с$^{-1}$},
				legend style={at={(0.02,0.97)}, anchor=north west},
				]
				\addlegendimage{/pgfplots/refstyle=plot:Ni}
				\addlegendentry{$N_i$ (левая ось)}

				\addplot[
					thick,
					mark=*,
					mark size=2pt,
				] coordinates {
					(4.50,0.10)
					(4.60,0.10)
					(4.70,0.90)
					(4.80,2.30)
					(4.90,1.70)
					(5.00,2.40)
					(5.10,1.50)
					(5.20,0.80)
					(5.30,0.10)
					(5.40,0.10)
				};
				\addlegendentry{$\rho_i$ (правая ось)}
			\end{axis}
		\end{tikzpicture}
		\caption{Число попаданий $N_i$ и плотность относительной частоты $\rho_i$ по интервалам ширины $\Delta t = 0.10$ с}
		\label{fig:N_rho}
	\end{figure}

	Так как $\rho_i = \dfrac{N_i}{N \, \Delta t}$, при фиксированных $N$ и $\Delta t$ плотность относительной частоты пропорциональна числу попаданий в интервал: обе зависимости имеют одинаковую форму и отличаются только масштабом.

	Максимум гистограммы находится на интервале $[4.95; 5.05)$ с. Для оценки координаты максимума возьмём середину этого интервала:
	\[
		t_{\text{гист}} \approx 5.00 \ \text{с}.
	\]
	Тогда
	\[
		t_{\text{гист}} - \bar t = 5.00 - 4.9429 = 0.0571 \ \text{с}.
	\]
	Это меньше и $\sigma = 0.1639$ с, и $\Delta t = 0.10$ с.

	Максимальная плотность нормального распределения:
	\[
		\rho_{\max} = \frac{1}{\sqrt{2 \pi} \sigma}
		\approx 2.435 \ \text{с}^{-1}.
	\]
	Максимальная высота экспериментальной гистограммы равна
	$2.40 \ \text{с}^{-1}$, то есть отличается от теоретической примерно на $1.4\%$.

	\begin{figure}[H]
		\centering
		\begin{tikzpicture}
			\begin{axis}[
				width=0.96\textwidth,
				height=8cm,
				xlabel={$t$, с},
				ylabel={Плотность относительной частоты $\rho$, с$^{-1}$},
				xmin=4.45, xmax=5.45,
				ymin=0, ymax=2.65,
				xtick={4.5,4.6,4.7,4.8,4.9,5.0,5.1,5.2,5.3,5.4},
				grid=major,
				legend style={at={(0.5,0.97)},anchor=north},
				]
				\addplot[
					ybar interval,
					draw=black,
					fill=gray!30,
				] coordinates {
					(4.45,0.10)
					(4.55,0.10)
					(4.65,0.90)
					(4.75,2.30)
					(4.85,1.70)
					(4.95,2.40)
					(5.05,1.50)
					(5.15,0.80)
					(5.25,0.10)
					(5.35,0.10)
					(5.45,0.10)
				};
				\addlegendentry{Гистограмма}

				\addplot[
					very thick,
					domain=4.45:5.45,
					samples=250,
				] {1/(sqrt(2*pi)*0.163869219844)*exp(-((x-4.9429)^2)/(2*(0.163869219844)^2))};
				\addlegendentry{Нормальное распределение}

				\addplot[dashed] coordinates {(4.9429,0) (4.9429,2.434516)};
			\end{axis}
		\end{tikzpicture}
		\caption{Гистограмма плотности относительной частоты и нормальное распределение}
		\label{fig:hist}
	\end{figure}

	\subsection{Сравнение гистограммы и нормального распределения}
	
	Гистограмма имеет один выраженный максимум и в целом близка к колоколообразной форме. Её максимум находится около $5.00$ с, тогда как максимум нормального распределения находится при $t = \bar t = 4.943$ с. Разность равна $0.057$ с и меньше $\sigma$ и ширины выбранного интервала $\Delta t$.

	Максимум гистограммы $2.40$ с$^{-1}$ близок к теоретическому $\rho_{\max} = 2.435$ с$^{-1}$. В целом форма гистограммы согласуется с нормальным распределением с параметрами $\bar t$ и $\sigma$.

	\section{Проверка интервалов $\bar t \pm \sigma$, $\bar t \pm 2 \sigma$, $\bar t \pm 3 \sigma$}

	Для нормального распределения теоретические вероятности попадания результата в интервалы $\bar t \pm \sigma$, $\bar t \pm 2 \sigma$ и $\bar t \pm 3 \sigma$ примерно равны $0.68$, $0.95$ и $0.997$ соответственно.

	\begin{table}[H]
		\centering
		\begin{xltabular}{\textwidth}{@{}X X X X X X@{}}
			\toprule
			Интервал & Нижняя граница, с & Верхняя граница, с & $N_{12}$ & $S_{\text{гист}} = N_{12} / N$ & $S_{\text{Гаусса}}$\\
			\midrule
			$\bar t \pm \sigma$   & 4.779 & 5.107 & 69  & 0.69 & 0.683 \\
			$\bar t \pm 2 \sigma$ & 4.615 & 5.271 & 96  & 0.96 & 0.954 \\
			$\bar t \pm 3 \sigma$ & 4.451 & 5.435 & 100 & 1.00 & 0.997 \\
			\bottomrule
		\end{xltabular}
		\caption{Доля результатов в интервалах $\bar t \pm k \sigma$: эксперимент и распределение Гаусса}
		\label{tab:sigma}
	\end{table}

	Полученные частоты $0.69$, $0.96$ и $1.00$ близки к теоретическим значениям $0.683$, $0.954$ и $0.997$. Это дополнительно согласуется с приближённым нормальным характером распределения выборки.

	\section{Случайная и полная погрешности}

	При $N = 100$ имеем $99$ степеней свободы. Для трёх требуемых доверительных вероятностей:

	\begin{table}[H]
		\centering
		\small
		\begin{xltabular}{\textwidth}{@{}X X X X X X X@{}}
			\toprule
			$\alpha$ & $t_{\alpha, N}$ & $\Delta t_{\text{сл}}$, с & $\varepsilon_{\text{сл}}$, \% &
			$\varepsilon_{\text{пр}}$, \% & $\varepsilon_{\text{полн}}$, \% & $\Delta t_{\text{полн}}$, с\\
			\midrule
			0.70 & 1.042 & 0.0171 & 0.35 & 0.20 & 0.40 & 0.020 \\
			0.90 & 1.660 & 0.0272 & 0.55 & 0.20 & 0.59 & 0.029 \\
			0.99 & 2.626 & 0.0430 & 0.87 & 0.20 & 0.89 & 0.044 \\
			\bottomrule
		\end{xltabular}
		\caption{Случайная, приборная и полная погрешности измерения}
		\label{tab:error}
	\end{table}

	Например, для $\alpha = 0.90$:
	\[
		\Delta t_{\text{сл}}
		= 1.660 \cdot 0.016387
		\approx 0.0272 \ \text{с},
	\]
	\[
		\varepsilon_{\text{сл}} = \frac{0.0272}{4.943} \approx 0.55 \, \%, \qquad
		\varepsilon_{\text{полн}} = \sqrt{0.55^2 + 0.20^2} \, \% \approx 0.59 \, \%,
	\]
	\[
		\Delta t_{\text{полн}} = 0.0059 \cdot 4.943 \approx 0.029 \ \text{с}.
	\]

	\section{Окончательные результаты}

	С учётом случайной и приборной погрешностей получаем:
	\[
		t = (4.943 \pm 0.020) \ \text{с}, \qquad \alpha = 0.70,
	\]
	\[
		t = (4.943 \pm 0.029) \ \text{с}, \qquad \alpha = 0.90,
	\]
	\[
		t = (4.943 \pm 0.044) \ \text{с}, \qquad \alpha = 0.99.
	\]

	\section{Вывод}

	\begin{enumerate}
		\item В результате $100$ измерений заданного промежутка времени получено среднее значение $\bar t = 4.943$ с при среднеквадратичном отклонении $\sigma \approx 0.164$ с и среднеквадратичной погрешности среднего $\sigma_{\bar t} \approx 0.016$ с. Окончательный результат с учётом случайной и приборной погрешностей: $t = (4.943 \pm 0.029)$ с при доверительной вероятности $\alpha = 0.90$.

		\item Полученная гистограмма плотности относительной частоты имеет один выраженный максимум и по форме в целом близка к колоколообразной. Её максимум равен $2.40$ с$^{-1}$ и отличается от теоретического значения $\rho_{\max} = 2.435$ с$^{-1}$ примерно на $1.4\%$. Положение максимума гистограммы (около $5.00$ с) отличается от $\bar t = 4.943$ с на $0.057$ с, что меньше как $\sigma$, так и ширины интервала $\Delta t = 0.10$ с.

		\item Относительные частоты попадания результатов в стандартные интервалы $\bar t \pm \sigma$, $\bar t \pm 2 \sigma$ и $\bar t \pm 3 \sigma$ составили соответственно $0.69$, $0.96$ и $1.00$, что близко к теоретическим значениям $0.683$, $0.954$ и $0.997$ для нормального распределения. Это, вместе с формой гистограммы, подтверждает приближённо нормальный характер распределения выборки.

		\item Измеренное среднее значение отличается от назначенного промежутка $5$ с на $5 - 4.943 = 0.057$ с. Эта разность превышает полную погрешность $\Delta t_{\text{полн}} \approx 0.029$ с при $\alpha = 0.90$, то есть носит систематический характер. Наиболее вероятная причина~--- время реакции участников измерения при ручном запуске и остановке таймера: команда <<стоп>> подаётся и обрабатывается с некоторой задержкой, из-за чего второй таймер, как правило, останавливается раньше, чем истекает ровно $5$ с.

		\item Ширина доверительного интервала существенно зависит от выбранной доверительной вероятности: при переходе от $\alpha = 0.70$ к $\alpha = 0.99$ полная погрешность увеличивается примерно в $2.2$ раза (с $0.020$ с до $0.044$ с). При этом уже при $\alpha = 0.70$ приборная погрешность $\Delta t_{\text{пр}} = 0.01$ с заметно влияет на результат, а при $\alpha = 0.99$ основной вклад вносит случайная составляющая. Это показывает, что при данном числе измерений случайная погрешность доминирует над приборной, и дальнейшее повышение точности разумно достигать увеличением числа измерений $N$.
	\end{enumerate}

%	\section{Контрольные вопросы}
%	
%	\begin{enumerate}
%		\item $\bar t$, $\sigma_{\bar t}$ и $\sigma$ имеют размерность времени; в данной работе они выражаются в секундах.
%		\item Дисперсия нормального распределения равна $D = \sigma^2$. В данной выборке $D = 0.026853 \ \text{с}^2$.
%		\item Коэффициент Стьюдента зависит от доверительной вероятности и числа степеней свободы. Он используется для расчёта случайной погрешности среднего значения: $\Delta t_{\text{сл}} = t_{\alpha, N} \sigma_{\bar t}$.
%	\end{enumerate}

\end{document}